\documentclass[11pt,a4paper]{tesis}

\usepackage{amsmath}
\usepackage{amssymb}
\usepackage{txfonts}
\usepackage{graphicx}
\usepackage[utf8]{inputenc}
\usepackage[spanish]{babel}
\usepackage[left=3cm,right=3cm,bottom=3.5cm,top=3.5cm]{geometry}

\newcommand{\until}{\text{ }\mathcal{U}\text{ }}
\newcommand{\release}{\text{ }\mathcal{R}\text{ }}
\newcommand{\since}{\text{ }\mathcal{S}\text{ }}
\newcommand{\trigger}{\text{ }\mathcal{T}\text{ }}
\newcommand{\weaku}{\text{ }\mathcal{W}\text{ }}

\newcommand{\next}{\medcirc}
\newcommand{\prev}{\medbullet}
\newcommand{\weakprev}{\tilde \medbullet}

\newcommand{\F}{\lozenge}
\renewcommand{\O}{\blacklozenge}

\newcommand{\G}{\square}
\renewcommand{\H}{\blacksquare}

\begin{document}

%%%% CARATULA

\def\autor{Joaquín Laks, Leonel Braginski, Tobias Moraut}
\def\tituloTesis{Volver al futuro: Incorporación de operadores de pasado al lenguaje lógico de MTSA}
\def\runtitulo{Volver al futuro: Incorporación de operadores de pasado al lenguaje lógico de MTSA}
\def\director{Sebastián Uchitel}
\def\lugar{Buenos Aires, 2026}
\newcommand{\HRule}{\rule{\linewidth}{0.2mm}}
%
\thispagestyle{empty}

\begin{center}\leavevmode

\vspace{-2cm}

\begin{tabular}{l}
\includegraphics[width=2.6cm]{logofcen.pdf}
\end{tabular}


{\large \sc Universidad de Buenos Aires

Facultad de Ciencias Exactas y Naturales

Departamento de Computaci\'on}

\vspace{6.0cm}

%\vspace{3.0cm}
%{
%\Large \color{red}
%\begin{tabular}{|p{2cm}cp{2cm}|}
%\hline
%& Pre-Final Version: \today &\\
%\hline
%\end{tabular}
%}
%\vspace{2.5cm}

\begin{huge}
\textbf{\tituloTesis}
\end{huge}

\vspace{2cm}

{\large Tesis de Licenciatura en Ciencias de la Computaci\'on}

\vspace{2cm}

{\Large \autor}

\end{center}

\vfill

{\large

{Director: \director}


\vspace{.2cm}

\lugar
}

\newpage\thispagestyle{empty}


%%%% ABSTRACTS, AGRADECIMIENTOS Y DEDICATORIA
\frontmatter
\pagestyle{empty}
%\begin{center}
%\large \bf \runtitulo
%\end{center}
%\vspace{1cm}
\chapter*{\runtitulo}

\noindent .

\bigskip

\noindent\textbf{Palabras claves:} .


%\cleardoublepage
%\chapter*{Agradecimientos}

\noindent Lorem ipsum dolor sit amet, consectetur adipiscing elit. Fusce sapien ipsum, aliquet eget convallis at, adipiscing non odio. Donec porttitor tincidunt cursus. In tellus dui, varius sed scelerisque faucibus, sagittis non magna. Vestibulum ante ipsum primis in faucibus orci luctus et ultrices posuere cubilia Curae; Mauris et luctus justo. Class aptent taciti sociosqu ad litora torquent per conubia nostra, per inceptos himenaeos. Mauris sit amet purus massa, sed sodales justo. Mauris id mi sed orci porttitor dictum. Donec vitae mi non leo consectetur tempus vel et sapien. Curabitur enim quam, sollicitudin id iaculis id, congue euismod diam. Sed in eros nec urna lacinia porttitor ut vitae nulla. Ut mattis, erat et laoreet feugiat, lacus urna hendrerit nisi, at tincidunt dui justo at felis. Class aptent taciti sociosqu ad litora torquent per conubia nostra, per inceptos himenaeos. Ut iaculis euismod magna et consequat. Mauris eu augue in ipsum elementum dictum. Sed accumsan, velit vel vehicula dignissim, nibh tellus consequat metus, vel fringilla neque dolor in dolor. Aliquam ac justo ut lectus iaculis pharetra vitae sed turpis. Aliquam pulvinar lorem vel ipsum auctor et hendrerit nisl molestie. Donec id felis nec ante placerat vehicula. Sed lacus risus, aliquet vel facilisis eu, placerat vitae augue.
 % OPCIONAL: comentar si no se quiere

%\cleardoublepage
%\hfill \textit{A mi persona favorita.}
  % OPCIONAL: comentar si no se quiere

\cleardoublepage
\tableofcontents

\mainmatter
\pagestyle{headings}

%%%% ACA VA EL CONTENIDO DE LA TESIS
\chapter{1}

\section{El lenguaje pLTL}

La gramática de LTL implementada hasta el momento en MTSA es:

\[
	\begin{array}{|l l l}
		\pi^i \models \top \\
		\pi^i \models \neg \varphi & \iff \pi^i \not\models \psi \\
		\pi^i \models \varphi \wedge \psi & \iff \pi^i \models \varphi \text{ y } \pi^i \models \psi \\
		\pi^i \models \varphi \vee \psi & \iff \pi^i \models \varphi \text{ o } \pi^i \models \psi \\
		\pi^i \models \next \varphi & \iff \pi^{i+1} \models \varphi & \textit{(Next, Siguiente)}\\
		\pi^i \models \varphi \until \psi & \iff \exists k \geq i / \pi^k \models \psi \text{ y } \forall i \leq j < k, \pi^j \models \varphi & \textit{(Until, Hasta)}\\
		\pi^i \models \varphi \release \psi & \iff \forall k \geq i, \pi^k \models \psi \text{ o } \exists i \leq j < k / \pi^j \models \varphi & \textit{(Release, Libera)}\\
	\end{array}
\]

De estos símbolos se desprenden las definiciones:

\[
	\begin{array}{|l l l}
		\bot & \equiv \neg \top & \textit{(Falso)} \\
		\F \varphi & \equiv \top \until \varphi & \textit{(Eventualmente $\varphi$)} \\
		\G \varphi & \equiv \neg \F \neg \varphi & \textit{(Siempre $\varphi$)} \\
		\varphi \weaku \psi & \equiv \psi \release (\varphi \vee \psi) & \textit{(Until débil)}
	\end{array}
\]

Estos últimos símbolos son parte del lenguaje LTL al que tiene acceso el usuario, pero el parser los transforma inmediatamente a su equivalente con los operadores de la primera tabla.

\citet{pnueli93} utiliza los operadores $\top, \neg, \wedge, \next, \until$, y agrega los operadores de pasado:

\[
	\begin{array}{|l l l}
		\pi^i \models \prev \varphi & \iff i > 0 \text{ y } \pi^{i-1} \models \varphi & \textit{(Yesterday, Anterior)}\\
		\pi^i \models \varphi \since \psi & \iff \exists k,  0 \leq k \leq i / \pi^k \models \psi \text{ y } \forall k \leq j < i, \pi^j \models \varphi & \textit{(Since, Desde)}\\
	\end{array}
\]

De los cuales, de la misma forma que en futuro, por comodidad del usuario, agregamos las definiciones:

\[
	\begin{array}{|l l l}
		\O \varphi & \equiv \top \since \varphi & \textit{(Eventualmente, en el pasado, $\varphi$)} \\
		\H \varphi & \equiv \neg \O \neg \varphi & \textit{(Siempre, en el pasado, $\varphi$)} \\
	\end{array}
\]

Por simplicidad, MTSA mantiene siempre las fórmulas en forma normal negada, es decir que el $\neg$ solo puede estar antes de una proposición. Para poder mantener esta restricción, cuando se crea una fórmula con una negación, se aplica la transformación recursiva:

\[
	\begin{array}{l l l l}
	\neg X \varphi & \rightarrow & \tilde X \neg \varphi & \text{Si $X$ es un operador unario y $\tilde X$ su dual}\\
		\neg (\psi X \varphi) & \rightarrow & \neg \psi \tilde X \neg \varphi & \text{Si $X$ es un operador binario y $\tilde X$ su dual}\\
		\neg \varphi & \rightarrow & \neg \varphi & \text{Si $\varphi$ es una proposición} \\
	\end{array}
\]

Queda claro entonces que es necesario que todo operador tenga un dual. Para todos los operadores de futuro ya presentes, existe en el lenguaje un dual ($\until$ y $\release$ son duales, y $\next$ es su propio dual). Pero \citet{pnueli93} no presenta operadores duales para los símbolos de pasado que introduce.

Para el $\since$, agregamos al lenguaje su dual $\trigger$, definido como:

\[
	\begin{array}{|l l l}
		\pi^i \models \varphi \trigger \psi & \iff \forall 0 \leq k \leq i, \pi^k \models \psi \text{ o } \exists k < j \leq i / \pi^j \models \varphi & \textit{(Trigger, Gatilla)}
	\end{array}
\]

A diferencia de su equivalente de futuro, no es verdad que $\prev$ es su propio dual. Esto se debe a la condición agregada $i>0$ que hace que $\pi^0 \not \models \prev \varphi$ para cualuier $\pi$ y $\varphi$. Entonces no es cierto que $\neg \prev \varphi \equiv \prev \neg \varphi$, porque en el momento 0, $\pi^0 \models \neg \prev \varphi$ pero $\pi^0 \not \models \prev \neg \varphi$. Por eso es necesario agregar el operador $\weakprev$, con semántica:


\[
	\begin{array}{|l l l}
		\pi^i \models \weakprev \varphi & \iff i = 0 \text{ o } \pi^{i-1} \models \varphi & \textit{(Yesterday débil)}
	\end{array}
\]

Es interesante notar que, evaluados en cualquier momento que no es el inicial, $\prev$ y $\weakprev$ son equivalentes.

Agregando $\trigger$ y $\weakprev$ a los operadores usados en \citet{pnueli93} y MTSA, obtenemos un lenguaje de lógica temporal con operadores de pasado y futuro que podemos mantener en forma normal negada.





Esta tesis habla mucho sobre logica temporal, es muy util. Por lo tanto propongo hacer un paper sobre esto. Waking up in the morning feeling like P.Diddy

%%%% BIBLIOGRAFIA
\backmatter
%\bibliography{tesis}

\end{document}
